% Contexto e instrucciones literales de las páginas 1 y 2 de EXA-0046.
% Fuente: Selectividad/sel_2014_matematicas_aplicadas/Matematicas_CCSS_JUN.pdf
% SHA-256: df7375f8183eafb6f45e9c65083e68d528ffabe07d66717b0ccf0f020c5b064f
% No se reproduce el logotipo ni el marco de la cabecera.
\section*{Contexto del examen original}
\textbf{UNIVERSIDADES DE ANDALUCÍA}\par
\textbf{PRUEBA DE ACCESO A LA UNIVERSIDAD}\par
CURSO 2013-2014\par
\textbf{MATEMÁTICAS APLICADAS A LAS CIENCIAS SOCIALES II}

\textbf{Instrucciones:}
\begin{enumerate}[label=\textbf{\alph*)}]
  \item Duración: 1 hora y 30 minutos.
  \item Elija una de las dos opciones propuestas y conteste los ejercicios de la opción elegida.
  \item En cada ejercicio, parte o apartado se indica la puntuación máxima que le corresponde.
  \item Se permitirá el uso de calculadoras que no sean programables, gráficas ni con capacidad para almacenar o transmitir datos.
  \item Si obtiene resultados directamente con la calculadora, explique con detalle los pasos necesarios para su obtención sin su ayuda. Justifique las respuestas.
\end{enumerate}
\noindent\textbf{Nota de la colección:} Se reúnen las dos opciones para consulta. La regla anterior corresponde al examen original. Curso 2013-2014 impreso. Convocatoria, tipo, reserva y modelo no acreditados; JUN sólo del nombre. Examen y criterios sin alias exactos en el inventario. Criterios concordantes; sin tabla específicamente asociada.
\par\medskip
